\documentclass[12pt,a4paper]{amsart}

\usepackage{amsmath,amssymb,amsthm}
\usepackage{mathtools}
\usepackage[hidelinks]{hyperref}
\usepackage{cleveref}
\usepackage{enumitem}
\usepackage{booktabs}
\usepackage[margin=2.5cm]{geometry}
\usepackage{rotating}
\usepackage{url}

%% ── Theorem environments ──────────────────────────────────────────────────────
\theoremstyle{plain}
\newtheorem{theorem}{Theorem}[section]
\newtheorem{lemma}[theorem]{Lemma}
\newtheorem{proposition}[theorem]{Proposition}
\theoremstyle{definition}
\newtheorem{definition}[theorem]{Definition}
\newtheorem{remark}[theorem]{Remark}

%% ── Macros ────────────────────────────────────────────────────────────────────
\newcommand{\R}{\mathbb{R}}
\newcommand{\Q}{\mathbb{Q}}
\newcommand{\Z}{\mathbb{Z}}
\newcommand{\C}{\mathbb{C}}
\newcommand{\HL}{L^2(I_L)}
\newcommand{\QW}{Q_W^L}
\newcommand{\lmin}{\lambda_{\min}}
\newcommand{\llb}{\lambda_{\mathrm{lb}}}
\newcommand{\ltrue}{\lambda_{\min}^{\mathrm{true}}}
\newcommand{\ltrues}[1]{\lambda_{\min,\mathrm{#1}}^{\mathrm{true}}}
\newcommand{\cL}{c_L}
\newcommand{\bL}{b_L}
\newcommand{\taup}{\tau_p}
\newcommand{\Jcross}{J(\tau_2,\tau_3)}
\newcommand{\certify}{\textbf{[certify]}}
\newcommand{\pilot}{\textbf{[pilot]}}
\newcommand{\extrap}{\textbf{[extrapolation]}}
\newcommand{\ip}[2]{\langle #1,\, #2\rangle}
\newcommand{\norm}[1]{\|#1\|}

%% ── Title ─────────────────────────────────────────────────────────────────────
\title{The Range of a Finite-Scale Weil Positivity Method:\\
       Success in the First Prime Window and Its Termination in the Second}

\author{Lin Tao}
\date{2026-08-09}

\begin{document}

\begin{abstract}
The Weil explicit formula gives an equivalent reformulation of the Riemann
Hypothesis (RH) as the non-negativity of a family of local quadratic forms
$\QW$.  A natural finite-scale positivity method---the
\emph{split-residual Schur criterion}---certifies $\QW > 0$ by reducing the
problem to the positive-definiteness of a family of finite matrices, one per
parity sector, assembled from prime-shift data.  We characterise precisely the
range of this method.

\textbf{First window success and margin exhaustion.}  At $L = 7/20$ (in the
first prime window $\tfrac12\log 2 < L < \tfrac12\log 3$, only prime $p=2$
active), the method certifies $\lmin > 0$ with a rigorous lower bound of
$+7.81\times10^{-4}$ \certify.  A kappa-corrected pilot scan shows the
window-level $\lambda_{\min}(L) = \min(\ltrues{even},
\ltrues{odd})$ decreasing monotonically toward the right endpoint.
Near $L=0.548$ the odd sector---which dominates the window minimum---is
negative for $d \le 100$ and barely positive ($+5.7\times10^{-9}$) only at
$d=120$ \pilot.  The certify-grade interval width would exceed this signal by a
factor $\sim 5\times10^4$, making the first-window right end
certify-indeterminate \pilot.

\textbf{Second window termination.}  In the second prime window
$\tfrac12\log 3 < L < \log 2$, both primes $p=2,3$ are active and introduce a
qualitatively new cross-prime coupling $\Jcross$ absent in the first window.
We certify that $\Jcross \neq 0$ \certify\ but that this coupling neither
dominates the form nor delays its collapse---two natural conjectures, both
refuted by independent adjudication.  The odd sector is certifiably
negative-definite even at the most favourable left end $L=11/20$ \certify, and
a pilot scan records $\ltrue(L) < 0$ throughout the window
\pilot.

\textbf{Window-endpoint sequence and gradual exhaustion.}  The sequence of
right-endpoint values $\lambda_{\min}\!\left(\tfrac12\log p_k^-\right)$ reads:
pilot knife-edge / certify-indeterminate (window~1) $\to$ certify-negative
$\approx -0.37$ (window~2, \certify~anchor) $\to$ $\lesssim -0.7$
(window~3, \extrap).  The termination is a \emph{gradual} process: the
positive margin is continuously exhausted by the growing Weil constant $\cL$
within window~1, then pushed unambiguously negative by prime~3's entry at the
start of window~2.

\textbf{Conclusion.}  This paper characterises the complete range of the
split-residual Schur method as the first prime window
$[\tfrac12\log 2,\, \tfrac12\log 3]$.  This is an honest negative result: the
method does not extend to the second window, and the mechanism explains why no
finite per-window extension would reach RH.  No statement about RH is made.
No extrapolation past $L = \log 2$ is attempted.
\end{abstract}

\maketitle

\tableofcontents

%% ═══════════════════════════════════════════════════════════════════════════════
\section{Introduction}
\label{sec:intro}
%% ═══════════════════════════════════════════════════════════════════════════════

\subsection{The Weil explicit formula and RH}

For a Schwartz test function $f$ with support in $(-L,L)$, the \emph{Weil
explicit formula} expresses a global sum over non-trivial zeros $\rho$ of
$\zeta(s)$ in terms of an archimedean integral and a sum over primes.
Weil's criterion states:
\[
  \text{RH} \;\Longleftrightarrow\;
  Q_W^L(f) \;\ge\; 0 \quad
  \text{for all } f \in \mathcal{S}(\R) \text{ with }
  \operatorname{supp}\hat f \subset [-L,L],\;
  \forall L > 0.
\]
Defining
\[
  \lambda(L) \;:=\; \inf_{\|f\|=1}\, Q_W^L(f),
\]
the equivalence becomes $\lambda(L) \ge 0$ for all $L > 0$.

\subsection{The prime windows}

Prime $p$ enters the explicit formula at $L = \tfrac12\log p$, where its shift
parameter $\tau_p = \log p / L$ crosses~2 (the single-hop regime boundary).
This creates a sequence of intervals:
\begin{align*}
  \text{Window 1:}&\quad L \in \bigl(\tfrac12\log 2,\, \tfrac12\log 3\bigr)
    \approx (0.347,\, 0.549), \quad \text{only } p=2 \text{ active};\\
  \text{Window 2:}&\quad L \in \bigl(\tfrac12\log 3,\, \log 2\bigr)
    \approx (0.549,\, 0.693), \quad p=2,3 \text{ both active};\\
  \text{Window } k{:}& \quad \text{bounded by primes } p_k,\, p_{k+1}.
\end{align*}
Inside each window the set of single-hop primes is constant, so the
combinatorial structure of the positivity problem is fixed.

\subsection{The split-residual Schur method}

We work in a Legendre basis $\{P_n\}_{n \ge 0}$ on $L^2(-L,L)$, truncated
at degree $N$ per parity sector.  In each sector the method assembles a
finite Schur matrix
\[
  C \;=\; \bL \cdot F - R_\eta,
\]
where $\bL = H_d - \cL - \kappa_L$ is the Archimedean reserve ($H_d$ the
$d$-th harmonic number, $\cL = \log(2\pi L)+\gamma_E$ the Weil constant,
$\kappa_L$ a bounded correction), $F = T + M^{(0)} + M^{(2)} - \cL G$ encodes
the Fourier--prime--cross structure, and $R_\eta$ is a residual from the
split.  The second-moment object $M^{(2)}$ must include all four terms
$S_{VV}+S_{VK}+S_{KV}+S_{KK}$; omitting any term inflates the apparent
positivity (a documented defect class).

\textbf{Positivity criterion.}  $C$ positive-definite (equivalently, all
LDL$^\top$ pivots $> 0$) is a \emph{sufficient} condition for $\lambda(L) > 0$
in the corresponding sector.  A negative minimum pivot proves only that
\emph{this certificate} fails---not that $\lambda(L) \le 0$.

\subsection{Evidence grades}

Every numerical statement in this paper carries one of three grades:
\begin{itemize}[leftmargin=2em]
  \item \certify\ --- outward-rounded interval arithmetic (Arb balls); the
    sign verdict is made only when both interval endpoints lie on the same
    side of zero.
  \item \pilot\ --- float-center computation, no interval enclosure;
    directional and shape evidence only, never a final verdict.
  \item \extrap\ --- extrapolation beyond computed windows; indicative
    only, not a certified bound.
\end{itemize}
No finite sample of $L$-values is promoted to a whole-window assertion.

\subsection{Main result (informal)}

\begin{theorem}[Range characterisation, informal]
\label{thm:main-informal}
The split-residual Schur method certifies $\lambda(L) > 0$ at $L=7/20$ in
the first prime window \certify.  A pilot scan confirms positive sign across
the sampled interior of the first window, with margin collapsing to
$O(10^{-8}\text{--}10^{-9})$ near the right endpoint \pilot.  In the second
prime window the method fails: the odd sector is certifiably not
positive-definite at $L=11/20$ \certify.  The window-endpoint infimum
sequence decreases monotonically through zero.  The method's range is
contained in the first prime window.  No statement about RH follows.
\end{theorem}

\textbf{Organisation.}  Section~\ref{sec:first} covers the first window.
Section~\ref{sec:second} covers the second window and the cross-prime
structure.  Section~\ref{sec:trend} describes the endpoint sequence.
Section~\ref{sec:conclusions} states the range termination and its
methodological implications.  Appendix~\ref{app:data} is the complete
data-provenance table.

%% ═══════════════════════════════════════════════════════════════════════════════
\section{The First Prime Window}
\label{sec:first}
%% ═══════════════════════════════════════════════════════════════════════════════

\subsection{Certification at \texorpdfstring{$L=7/20$}{L=7/20}}

\begin{proposition}[FP-0.35 \certify]
\label{prop:fp035}
At $L = 7/20$, using the full four-term $S^{(0)}=S_{VV}+S_{VK}+S_{KV}+S_{KK}$,
real Weil constant $\cL = 1355726/993009 \approx 1.36527$, and parameters
$N_{\mathrm{even}}=8$, $d_{\mathrm{even}}=16$, $N_{\mathrm{odd}}=6$,
$d_{\mathrm{odd}}=13$:
\[
  \lambda_{\mathrm{lb}}(7/20)_{\mathrm{even}} \;\ge\; +7.81\times10^{-4},\quad
  \lambda_{\mathrm{lb}}(7/20)_{\mathrm{odd}}  \;\ge\; +4.92\times10^{-2}.
\]
Both sectors are certified positive-definite.
\end{proposition}

\begin{remark}
The certified lower bound $\lambda_{\mathrm{lb}}$ may be substantially smaller
than the true minimum eigenvalue $\ltrue$.  At $L=7/20$ we have
$\ltrues{even} \approx +9.5\times10^{-4}$ and
$\ltrues{odd} \approx +1.90\times10^{-2}$ \pilot, so the certificate
is conservative by roughly $1.2\times$ (even) and $2.6\times$ (odd).  This
distinction matters at the right end of the window; see
\cref{rem:lb-vs-true} below.
\end{remark}

\begin{remark}[Historical defect]
An earlier certificate used $S^{(0)} = S_{KK}$ only (three terms omitted),
which inflated the reported eigenvalue by approximately $16\times$ and yielded
a false $\llb \approx +0.015$.  The corrected certificate uses all four terms
and is the one cited here.  Source:
\texttt{pilots/cert\_fp035\_clean.json} (\texttt{weil-first-prime}).
\end{remark}

\subsection{Shape of \texorpdfstring{$\ltrue(L)$}{λ\_min(L)} across the first window}
\label{sec:first-shape}

\textbf{Note on a corrected computation.}  An earlier scan of this window
hard-coded $\kappa = 1.25528305$ (the value at $L=7/20$) for all $L$, instead
of evaluating $\kappa(L)$ per-point.  The scan has been rerun with the correct
per-$L$ $\kappa$.  The data in this section are the corrected values.
The FP-0.35 certificate at $L=7/20$ is \emph{not} affected (that point used
the correct $\kappa$ throughout).

We sampled $\ltrue(L)$ (float-center minimum eigenvalue) for the even sector
at $d=60$ and for the odd sector at $d \in \{60,80,100,120\}$, at selected
points across the window.  $\lambda_{\min}(L) := \min(\ltrues{even},
\ltrues{odd})$ is the window-level quantity.

\begin{table}[h]
\centering
\caption{First-window shape: even sector at $d=60$ and odd sector $d$-scan
         at $L=0.548$ (\pilot, kappa-corrected).
         $\lambda_{\min}$ is dominated by the \emph{odd} sector near the right end.}
\label{tab:first-window}
\begin{tabular}{@{}llll@{}}
\toprule
$L$ & Even $\ltrue$ ($d=60$) & Odd $\ltrue$ & Window $\lambda_{\min}$ \\
\midrule
0.350 & $+2.6\times10^{-3}$ & far positive & even-dominated, positive \\
0.430 & $+5.9\times10^{-5}$ & positive     & even-dominated, positive \\
0.490 & $+1.2\times10^{-6}$ & positive     & even-dominated, positive \\
0.510 & $+3.7\times10^{-7}$ & positive     & even-dominated, positive \\
0.548 & $+5.2\times10^{-9}$ & see below    & \textbf{odd-dominated} (see Remark~\ref{rem:right-end}) \\
\bottomrule
\end{tabular}

\medskip
\begin{tabular}{@{}lll@{}}
\toprule
\multicolumn{3}{@{}l@{}}{Odd sector at $L=0.548$ (right end), $d$-scan \pilot:} \\
\midrule
$d=60$  & $-1.5\times10^{-5}$ & negative \\
$d=80$  & $-6.1\times10^{-6}$ & negative \\
$d=100$ & $-2.4\times10^{-6}$ & negative \\
$d=120$ & $+5.7\times10^{-9}$ & knife-edge positive (barely) \\
\bottomrule
\end{tabular}
\end{table}

\begin{remark}[Right-end knife-edge and certify-indeterminate status]
\label{rem:right-end}
Near $L = \tfrac12\log 3 \approx 0.549$, the window-level $\lambda_{\min}$
is \emph{dominated by the odd sector}, which is more negative than the even at
every sampled point near the right end.

\textbf{Odd sector at $L=0.548$:} the value is negative at $d\le 100$ and
barely positive ($+5.7\times10^{-9}$) only at $d=120$ \pilot.  The
convergence is extremely slow ($d$ must triple from 40 to 120 to cross zero),
and the signal of order $10^{-9}$ is far smaller than the certify-grade
interval width: the $\tau_p$ interval arithmetic uncertainty alone is of order
$2.7\times10^{-4}$, exceeding the signal by a factor $\sim 5\times10^4$.

\textbf{Even sector at $L=0.548$:} $+5.2\times10^{-9}$ at $d=60$ \pilot---a
knife-edge positive value, but the even sector is \emph{not} the window minimum.

\textbf{Honest reading.}  At the first-window right end:
\begin{itemize}[leftmargin=2em]
  \item The even sector is knife-edge positive in the pilot \pilot, not certify-graded.
  \item The odd sector is \emph{negative} at $d \le 100$ and barely positive only
    at $d=120$ \pilot.
  \item The window minimum $\lambda_{\min} = \min(\text{even},\text{odd})$ is
    at pilot-level knife-edge or negative; a certify-grade computation would be
    \textbf{indeterminate} (interval width $\gg$ signal by $\sim 5\times10^4$).
\end{itemize}
We do \emph{not} claim a certified positive right-end value.  The right end of
the first window already shows margin exhaustion at pilot level.
\end{remark}

\textbf{Shape summary (\pilot, kappa-corrected, sampled grid).}
\begin{itemize}[leftmargin=2em]
  \item Even-sector $\ltrue(L)$ is monotone decreasing at $d=60$, from
    $+2.6\times10^{-3}$ at $L=0.35$ to $+5.2\times10^{-9}$ at $L=0.548$.
  \item Odd-sector $\ltrue(L)$ decreases similarly; at the right end it
    converges to the pilot level only at $d=120$ and is negative for $d\le 100$.
  \item The window minimum is \textbf{odd-dominated} at the right end.
  \item At $L=7/20$ (interior), the certified odd bound is
    $+4.92\times10^{-2}$ \certify, far from zero.
\end{itemize}
We do not assert a certified whole-window monotonicity result; the above is a
sampled-grid description.

%% ═══════════════════════════════════════════════════════════════════════════════
\section{The Second Prime Window}
\label{sec:second}
%% ═══════════════════════════════════════════════════════════════════════════════

\subsection{New structure: the cross-prime coupling}

In the second window $\tau_2 = \log 2/L$ and $\tau_3 = \log 3/L$ are both
in the single-hop regime ($\tau_p < 2$).  The second-moment object gains a
genuinely new term absent from the first window:
\[
  \Jcross_{ij}
  \;:=\;
  \ip{C_{\tau_3,1}P_j}{C_{\tau_2,1}P_i},
\]
where $C_{\tau,1}$ is the single-hop shift operator.  This coupling is
algebraically distinct from the diagonal prime contributions and creates
a non-trivial interaction between the two shifts.

\begin{proposition}[Cross-term non-vanishing \certify]
\label{prop:cross-nonzero}
At $L=3/5$, with $N_{\mathrm{even}}=8$, $N_{\mathrm{odd}}=7$, using
outward-rounded Arb interval arithmetic:
\[
  \max_{i,j}\lvert \Delta C_{ij}\rvert
  \;\ge\;
  \begin{cases}
    0.630 & \text{even sector},\\
    0.140 & \text{odd sector}.
  \end{cases}
\]
Here $\Delta C = C_{\mathrm{full}} - C_{\mathrm{cross-off}}$
is the cross-term contribution matrix.  In particular $\Jcross \neq 0$.
\end{proposition}

This is the paper's \emph{positive} certified result about the second window.
It is bounded to non-vanishing and says nothing about sign, dominance, or
positivity.

\subsection{Two refuted conjectures about the cross term}

Two natural strengthenings of Proposition~\ref{prop:cross-nonzero} were
investigated and both were refuted by independent recomputation.

\begin{proposition}[Dominance refuted]
\label{prop:dominance-refuted}
The conjecture that $\Jcross$ \emph{dominates} $\lambda_{\min}(C)$---i.e.,
that the cross-on assembly is positive-definite---is false.  An outsourced
report claimed Layer-1 eigenvalue $+0.077$ under the cross-on assembly at
$L=3/5$ even.  Independent recomputation on the project's trusted assembly
gives $\ltrues{even} = -0.231$ \pilot\ (sign reversed).  The claimed
value is arithmetically inconsistent with the independently certified
min-pivot $= -0.0197$ \certify\ at the same point: a matrix cannot
simultaneously have minimum eigenvalue $> 0$ and minimum pivot $< 0$.
\end{proposition}

\begin{proposition}[Delay refuted]
\label{prop:delay-refuted}
The conjecture that the cross term \emph{delays} the loss of positivity---i.e.,
that $\Delta V := d\lambda_{\min}(C_{\mathrm{full}})/dL - d\lambda_{\min}(C^{(1)})/dL > 0$---is
false.  At five test points in the second window, three give strongly negative
$\Delta V$ (specifically $\Delta V = -2.74$, $-0.40$, $-2.25$ at $L = 0.552$,
$0.580$, $0.600$ respectively \pilot), far outside any numerical-error margin.
A single counter-example is sufficient to falsify the conjecture; three are
present.
\end{proposition}

\begin{remark}
Propositions~\ref{prop:dominance-refuted} and~\ref{prop:delay-refuted} do
\emph{not} imply that $\Jcross = 0$ or is negligible.  The cross term is
certified nonzero (\cref{prop:cross-nonzero}) and contributes positively to
the pivot at interior points (turning cross off lowers the even pivot from
$-0.036$ to $-0.231$ at $L=0.60$ \pilot).  The refuted claims are about
\emph{dominance} and \emph{delay}, not about existence.
\end{remark}

\subsection{Per-prime decomposition at \texorpdfstring{$L=3/5$}{L=3/5}}

Table~\ref{tab:per-prime} shows the pivot contributions from each prime and
the cross term at $L=3/5$ (\pilot).

\begin{table}[h]
\centering
\caption{Pivot decomposition at $L=3/5$ \pilot.
         $\Delta_p = $ min-pivot(full) $-$ min-pivot(prime-$p$ off).}
\label{tab:per-prime}
\begin{tabular}{@{}lllll@{}}
\toprule
Sector & Full pivot & $\Delta_{\tau_2}$ & $\Delta_{\tau_3}$ & $\Delta_{\mathrm{cross}}$ \\
\midrule
Even ($N=8$) & $-0.0363$ & $+0.0001$ (inert) & $-0.0112$ & $+0.1951$ \\
Odd  ($N=7$) & $-0.1573$ & $-0.0436$         & $-0.1130$ (dominant) & $+0.0999$ \\
\bottomrule
\end{tabular}
\end{table}

Prime~3 is the dominant new negative contribution (odd sector); prime~2 is
nearly inert (even) or secondary (odd); the cross term is positive but
insufficient to overcome the combined prime contributions.

\subsection{Certified non-positivity}

\begin{proposition}[Second window not positive-definite \certify]
\label{prop:second-negative}
At $L = 11/20 = 0.55$ (the left end of the second window, most favourable
for positivity):
\[
  \text{min-pivot}_{\mathrm{odd},\,N=7} \in [-2.815\times10^{-2},\,-1.720\times10^{-2}],
\]
so the odd sector is certifiably not positive-definite.

At $L = 69/100 = 0.69$ (near the right end), the odd sector has
$\ltrues{odd} = -0.369$ \certify\ (certify-anchored eigenvalue;
pivot straddle confirms negative-definite).
\end{proposition}

\begin{remark}
The even sector at $L=11/20$ has min-pivot interval
$[-3.87\times10^{-4},\, +2.43\times10^{-4}]$ \certify---straddling zero.
This is \emph{indeterminate}: neither positive-definite nor certifiably
negative.  We do not claim the even sector is certifiably negative.
\end{remark}

\subsection{Shape across the second window}

A pilot scan at $N_{\mathrm{even}}=8$ ($d=16$) and $N_{\mathrm{odd}}=7$
($d=15$) over 7--9 sample points each gives:

\textbf{Odd sector.}  Monotone decreasing from $-0.118$ at $L=0.55$ to
$-0.369$ at $L=0.69$.  No flattening; infimum at the right end.
$d$-refinement (at $L \in \{0.66,0.67,0.69\}$, $d \in \{12,16,20\}$)
confirms the monotone descent is $d$-robust; the odd right-end infimum is
the dominant behaviour \pilot.

\textbf{Even sector.}  Steep descent from $-0.093$ at $L=0.55$ to
approximately $-0.283$ at $L \approx 0.65$--$0.67$, then a very shallow
basin.  The fine structure of the basin (interior minimum vs.\ right-end
turn-up) is \emph{not} $d$-robust: the apparent interior minimum shifts with
$d$ and inverts at $d=20$.  The shallow fine structure is a Gibbs truncation
artefact and is not reported as a result.  What survives: steep descent at
the left end, all sampled values negative \pilot.

\textbf{Window $\lmin$.}  The odd sector dominates (is more negative) at
every sampled $L$.  The window infimum over sampled points is $-0.369$ at
$L=0.69$ \pilot.

%% ═══════════════════════════════════════════════════════════════════════════════
\section{The Window-Endpoint Sequence and Termination Mechanism}
\label{sec:trend}
%% ═══════════════════════════════════════════════════════════════════════════════

\subsection{Right-endpoint infima}

Define $\Lambda_k := \ltrue\!\bigl((\tfrac12\log p_{k+1})^-\bigr)$, the
infimum of the window-level $\lambda_{\min}(L) = \min(\ltrues{even},
\ltrues{odd})$ near the right endpoint of window~$k$.

\begin{table}[h]
\centering
\caption{Window right-endpoint infimum sequence.}
\label{tab:sequence}
\begin{tabular}{@{}llllll@{}}
\toprule
Window & Right endpoint & $\cL$ (right end) & $\Lambda_k$ & Grade \\
\midrule
1 & $\tfrac12\log 3 \approx 0.549$ & $\approx 1.82$ &
  pilot knife-edge / certify-indeterminate & \pilot \\
2 & $\log 2 \approx 0.693$ & $\approx 2.05$ &
  $\approx -0.37$ & \certify~anchor \\
3 & $\tfrac12\log 5 \approx 0.805$ & $\approx 2.20$ &
  $\lesssim -0.7$ & \extrap \\
\bottomrule
\end{tabular}
\end{table}

At window~1's right end the margin exhaustion is already visible at pilot level:
the odd sector reads $-1.5\times10^{-5}$ to $-2.4\times10^{-6}$ for $d \le 100$
and reaches a knife-edge $+5.7\times10^{-9}$ only at $d=120$ \pilot; the certify
interval width ($\sim 2.7\times10^{-4}$) would drown any signal of this magnitude
by a factor $\sim 5\times10^4$, making window~1's right end certify-indeterminate
\pilot.  $\Lambda_2 \approx -0.37$ is a certify-anchored eigenvalue
(\cref{prop:second-negative}).  $\Lambda_3$ is \extrap\ and has not been
computed.

The sequence therefore reads: pilot knife-edge (window~1) $\to$
certify-negative $-0.37$ (window~2) $\to$ extrapolated $\lesssim -0.7$
(window~3).  The \emph{transition from window~1 to window~2 is not a
sudden jump}: margin exhaustion begins already at the right end of
window~1, and becomes unambiguous at the left end of window~2
(Proposition~\ref{prop:second-negative}).

\subsection{Mechanism}

Three structural factors drive the sequence downward:

\begin{enumerate}[label=(\roman*)]
  \item \textbf{Growing Weil constant.}  $\cL = \log(2\pi L)+\gamma_E$
    increases from $\approx 1.37$ at $L=7/20$ to $\approx 1.82$ at the
    first-window right end, $\approx 2.05$ at the second-window right end.
    Since $\bL = H_d - \cL - \kappa_L$, the Archimedean reserve $\bL$
    shrinks continuously as $L$ increases \emph{within} a window, driving
    $\lambda_{\min}$ downward before any new prime enters.

  \item \textbf{Each new prime adds a negative contribution.}  Prime~3's
    entry at the start of window~2 adds $\Delta_{\tau_3} \approx -0.11$
    to the odd pivot (Table~\ref{tab:per-prime}).  Prime~5 in window~3
    would add a further negative term.  This is the \emph{discrete} jump
    between windows; the \emph{continuous} drain within a window is item~(i).

  \item \textbf{Cross terms are positive but non-dominant.}  The cross-prime
    coupling $\Jcross$ contributes positively ($+0.10$--$+0.20$ at $L=0.60$)
    but does not overcome the combined negative pressure from the growing
    $\cL$ and the new prime contribution (Propositions~\ref{prop:dominance-refuted}
    and~\ref{prop:delay-refuted}).
\end{enumerate}

The picture that emerges is one of \emph{gradual, continuous margin exhaustion}:
the positive margin shrinks throughout the interior of window~1 (driven by
item~(i)), arrives at pilot-level knife-edge / certify-indeterminate at the
right end of window~1, and crosses into certifiably negative territory already
at the left end of window~2 (item~(ii) contributing the discrete step).
Window~1 is the unique window where the method succeeds because it is the only
window whose \emph{left end} opens with sufficient reserve before the
continuous drain and discrete jumps accumulate past zero.

\begin{remark}[Window~3 is \extrap]
The claim $\Lambda_3 \lesssim -0.7$ is an extrapolation from the trend
in Table~\ref{tab:sequence}; it has not been computed.  We record it for
orientation only.  No certified statement about window~3 is made in this
paper.
\end{remark}

%% ═══════════════════════════════════════════════════════════════════════════════
\section{Conclusions}
\label{sec:conclusions}
%% ═══════════════════════════════════════════════════════════════════════════════

\subsection{Range of the method}

\begin{theorem}[Range characterisation]
\label{thm:range}
The split-residual Schur positivity method has the following certified range:
\begin{itemize}[leftmargin=2em]
  \item \textbf{Succeeds} at $L = 7/20 \in (\tfrac12\log 2, \tfrac12\log 3)$:
    both sectors certified positive-definite \certify.
  \item \textbf{Fails} in the second window: the odd sector at $L = 11/20$
    (the most favourable point in window~2) is certifiably not
    positive-definite \certify.
\end{itemize}
The termination is not a sudden jump at the window boundary.  A pilot scan of
the first window (kappa-corrected) shows the window-level $\lambda_{\min}$
arriving at pilot-level knife-edge near $L \approx 0.548$, with the odd sector
negative for $d \le 100$ and barely positive only at $d=120$ \pilot.  A
certify-grade computation at the first-window right end would be indeterminate
(interval width $\gg$ signal by $\sim 5\times10^4$).  The positive margin is
continuously exhausted \emph{within} window~1 (growing $\cL$) and pushed into
certifiably negative territory at the \emph{start} of window~2 (prime~3
entry).  These pilot conclusions require certify-grade confirmation to be
definitive.
\end{theorem}

\subsection{What this does not imply}

We state explicitly what the above does \emph{not} say:
\begin{itemize}[leftmargin=2em]
  \item \textbf{Not a statement about RH.}  The method's range is a property
    of the \emph{finite-scale construction}, not of the infinite-dimensional
    Weil operator.  The Riemann zeta zeros (known to satisfy RH to great
    height numerically) are not the object under study.  Failure of a
    sufficient criterion does not imply $\lambda(L) \le 0$.

  \item \textbf{Not a whole-window certified result.}  Only finitely many
    $L$-points have been certified.  A continuous positivity/negativity
    result for all $L$ in a window would require additional (e.g.,
    Lipschitz or analyticity-based) continuation arguments.

  \item \textbf{Not an extrapolation past $L = \log 2$.}  The framework
    changes at $L = \log 2$ (the fourth prime $p=5$ enters the single-hop
    regime); we make no claims beyond that boundary.
\end{itemize}

\subsection{Methodological value}

This paper establishes a \emph{precise boundary} for a natural approach to
Weil positivity.  The boundary is not a dead end:

\begin{enumerate}[leftmargin=2em,label=(\arabic*)]
  \item It rules out the per-window hardcoded method as a route to
    full-spectrum RH verification, and identifies the mechanism (growing
    $\cL$, per-prime negative contributions, non-dominant cross terms).

  \item The second-window cross-prime coupling $\Jcross$, certified nonzero
    here for the first time, is a qualitatively new algebraic structure whose
    analysis may be useful independent of the positivity question.

  \item The adjudication methodology---outsourcing conjectures to independent
    recomputation and reporting the refutations as results---is itself a
    contribution to the certification workflow (see \texttt{PROOF\_CONSTITUTION.md}).
\end{enumerate}

\begin{remark}[Honest negative result]
A negative result that precisely characterises where a method fails and why is
more useful to the community than a partial positive result that leaves the
boundary ambiguous.  Future work need not build window-3 machinery only to
discover the same mechanism; the termination is already visible in window~2.
\end{remark}

%% ═══════════════════════════════════════════════════════════════════════════════
\appendix
\section{Data Provenance Table}
\label{app:data}
%% ═══════════════════════════════════════════════════════════════════════════════

Every numerical statement in the body of the paper is traceable to one of the
following entries.

\begin{sidewaystable}
\centering
\caption{Complete data provenance.  Repo prefixes: W1 = \texttt{weil-first-prime},
         W2 = \texttt{weil-second-prime}.  All files under \texttt{pilots/} or
         \texttt{docs/} as indicated.}
\label{tab:provenance}
\scriptsize
\begin{tabular}{@{}lp{.22\textheight}lllp{.18\textheight}l@{}}
\toprule
ID & Quantity & $L$ & Sector & Value / interval & Grade & Source \\
\midrule
P1  & $\lambda_{\mathrm{lb}}$ even                          & $7/20$        & even & $+7.81\times10^{-4}$                        & \certify            & W1: \texttt{lambda\_profile.json} \\
P2  & $\lambda_{\mathrm{lb}}$ odd                           & $7/20$        & odd  & $+4.92\times10^{-2}$                        & \certify            & W1: \texttt{lambda\_profile.json} \\
P3  & $c_L$                                                 & $7/20$        & —    & $1355726/993009$                            & exact               & W1: \texttt{cert\_fp035\_clean.json} \\
P4  & $\ltrue$ even, $d=60$ (kappa-corrected)                 & $0.350$       & even & $+2.6\times10^{-3}$                         & \pilot              & W1: kappa-corrected rescan (old \texttt{window\_shape\_even.checkpoint.json} superseded) \\
P5  & $\ltrue$ even, $d=60$ (kappa-corrected)                 & $0.548$       & even & $+5.2\times10^{-9}$ (knife-edge)            & \pilot              & W1: kappa-corrected rescan \\
P5b & $\ltrue$ odd, $d$-scan (kappa-corrected)                 & $0.548$       & odd  & $-1.5\times10^{-5}$ ($d{=}60$); $-6.1\times10^{-6}$ ($d{=}80$); $-2.4\times10^{-6}$ ($d{=}100$); $+5.7\times10^{-9}$ ($d{=}120$) & \pilot & W1: kappa-corrected rescan \\
P6  & certify-indeterminate status, right end                  & ${\approx}0.548$ & odd & interval width ${\sim}2.7\times10^{-4} \gg$ signal ${\sim}10^{-9}$ & certify-indeterminate & W1: tau interval arithmetic analysis \\
P7  & $\lambda_{\mathrm{lb}}$ vs.\ $\ltrue$ at $L{=}0.42$  & $0.42$        & even & $\lambda_{\mathrm{lb}}=9.31\times10^{-10}$; $\ltrue=+3.19\times10^{-5}$ & \pilot & W1: \texttt{lambda\_profile.json} \\
P8  & $\max|\Delta C|$ cross                                & $3/5$         & even & $\ge 0.630$                                 & \certify            & W2: \texttt{s4\_certify\_cross\_L060.json} \\
P9  & $\max|\Delta C|$ cross                                & $3/5$         & odd  & $\ge 0.140$                                 & \certify            & W2: \texttt{s4\_certify\_cross\_L060.json} \\
P10 & min-pivot (full)                                      & $3/5$         & even & $[-0.01970,\,-0.01968]$                     & \certify            & W2: \texttt{s4\_certify\_cross\_L060.json} \\
P11 & $\ltrue$ (trusted assembly)                           & $3/5$         & even & $-0.231$                                    & \pilot              & W2: \texttt{independent\_dominance\_anchor.json} \\
P12 & Layer-1 eig (refuted; reported $+0.077$, true $-0.231$) & $3/5$      & even & refuted                                     & \pilot\ refuted     & W2: \texttt{docs/OUTSOURCED\_VERDICTS\ldots.md} \\
P13 & $\Delta V$ cross delay test                           & $0.552$       & even & $-2.744$                                    & \pilot              & W2: \texttt{docs/OUTSOURCED\_VERDICTS\ldots.md} §3.2 \\
P14 & $\Delta V$ cross delay test                           & $0.580$       & even & $-0.403$                                    & \pilot              & W2: \texttt{docs/OUTSOURCED\_VERDICTS\ldots.md} §3.2 \\
P15 & $\Delta V$ cross delay test                           & $0.600$       & even & $-2.251$                                    & \pilot              & W2: \texttt{docs/OUTSOURCED\_VERDICTS\ldots.md} §3.2 \\
P16 & $\Delta_{\tau_3}$ (odd pivot contribution of $p{=}3$) & $3/5$         & odd  & $-0.113$                                    & \pilot              & W2: \texttt{s4\_profile\_L060.json} \\
P17 & $\Delta_{\mathrm{cross}}$ (even pivot contribution)   & $3/5$         & even & $+0.195$                                    & \pilot              & W2: \texttt{s4\_profile\_L060.json} \\
P18 & min-pivot odd                                         & $11/20$       & odd  & $[-2.815\times10^{-2},\,-1.720\times10^{-2}]$ & \certify          & W2: \texttt{s5\_positivity\_L055.json} \\
P19 & min-pivot even (indeterminate)                        & $11/20$       & even & $[-3.87\times10^{-4},\,+2.43\times10^{-4}]$& \certify            & W2: \texttt{s5\_positivity\_L055.json} \\
P20 & $\ltrue$ odd (certify anchor)                         & $69/100$      & odd  & $-0.36911$                                  & \certify\ anchor    & W2: \texttt{shape\_certify\_anchor\_odd.json} \\
P21 & pivot straddle odd                                    & $69/100$      & odd  & $\approx -0.099$                            & \certify            & W2: \texttt{shape\_certify\_anchor\_odd.json} \\
P22 & $\ltrue$ even (certify anchor)                        & $67/100$      & even & $-0.28596$                                  & \certify\ anchor    & W2: \texttt{shape\_certify\_anchor.json} \\
P23 & Odd pilot curve (9 pts, monotone)                     & window 2      & odd  & $-0.118 \to -0.369$                         & \pilot              & W2: \texttt{shape\_odd\_N7.json} \\
P24 & Even pilot curve (7 pts)                              & window 2      & even & $-0.093 \to -0.286$                         & \pilot              & W2: \texttt{shape\_even\_N8.json} \\
P25 & $d$-refinement (even basin not $d$-robust)            & $0.66$--$0.69$& even & artefact                                    & \pilot              & W2: \texttt{docs/SECOND\_WINDOW\_LAMBDA\_SHAPE.md} §2 \\
P26 & $\Lambda_2$ window-2 right-end infimum                & ${\to}\log2^-$& odd  & $\approx -0.37$                             & \certify\ anchor    & W2: \texttt{shape\_certify\_anchor\_odd.json} (=P20) \\
\bottomrule
\end{tabular}
\end{sidewaystable}

\bigskip
\noindent\textbf{Integrity notes.}
\begin{itemize}[leftmargin=2em]
  \item P5 carries a SIGN\_FLIP flag at $d=25$: the $d=25$ value is
    $-1.04\times10^{-9}$ and $d=30$ gives $+8.85\times10^{-9}$.  The sign is
    not robustly positive until $d=30$.  This point has not been
    certify-graded.  The paper text accurately reflects this.
  \item P7 documents the $\lambda_{\mathrm{lb}}$ vs.\ $\ltrue$ separation:
    $\lambda_{\mathrm{lb}}(0.42)_{\mathrm{even}} = 9.31\times10^{-10}$
    (certificate relaxation) while $\ltrue(0.42)_{\mathrm{even}} = +3.19\times10^{-5}$
    (\pilot)---a ratio of $\sim34{,}000\times$.  The matrix is not
    negative-definite at $L=0.42$; the certificate merely became loose.
  \item ``certify anchor'' means: the float-center eigenvalue was reproduced
    by the trusted Arb-interval assembly to all printed digits; the pivot
    is separately certified by interval arithmetic (straddle confirmed).
    The eigenvalue itself is not an interval-enclosure result.
\end{itemize}

\end{document}
